% !TeX program = xelatex
% UTF-8. Compile with XeLaTeX from this directory or the repository root.
% Define \IDMInNotebook before inputting into an existing XePersian document.
% Source: Google Drive Notebook/infants od DM mothers (original folder spelling).
% pic01.jpg: 20261016_213238706.jpg, 1SVAjZDdp_QPgP80ufB8KTHHFig5J4fOe
% pic02.jpg: 20261016_213245233.jpg, 1B9y3BRScf5EgGNziAwuRzyA2AofFLBS6
% pic03.jpg: 20261016_213250727.jpg, 1xfrQmHgCQVlCb90cg2dfK6MxBEVm1aAR
% pic04.jpg: 20261016_213259432.jpg, 1LxF_nt2iWMRIScKtlVn2j7uegelKZPHC
% pic05.jpg: 20261016_213304290.jpg, 1YXukgZvDK4s3PSBOPb0PP1Q39hdpiRDa
% pic06.jpg: 20261016_213310734.jpg, 1vCdnioK-zttFbBPw-nGMdkEgME73TzU5
% pic07.jpg: 20261016_213316070.jpg, 1zFFaD014erMnQZzN89uOwzJdQMZKUl63
% Arrow relationships are represented by nested statements; source claims are
% transcribed, not clinically revised. Uncertain handwriting is marked explicitly.
\ifdefined\IDMInNotebook
  \let\IDMfinish\relax
\else
  \def\IDMfinish{\end{document}}
  \documentclass[12pt,a4paper]{article}
  \usepackage[margin=20mm]{geometry}
  \usepackage{amsmath,amssymb}
  \usepackage{graphicx}
  \usepackage{enumitem}
  \usepackage{xepersian}
  \IfFontExistsTF{Amiri}{\settextfont{Amiri}}{\settextfont{DejaVu Sans}}
  \setlatintextfont{DejaVu Serif}
  \setlength{\parindent}{0pt}
  \setlength{\parskip}{4pt}
  \setlength{\emergencystretch}{3em}
  \setlist[itemize]{nosep,topsep=3pt}
  \begin{document}
\fi
\providecommand{\IDMmath}[1]{\lr{\(#1\)}}
\providecommand{\IDMunclear}[1]{\textbf{[خوانش نامطمئن: #1]}}
\providecommand{\IDMpage}[1]{%
  \section*{تصویر #1 ـ ۱۴۷}
  {\footnotesize\lr{Subject:}\quad\lr{Date: / /}\par}}

\begin{center}
{\Large\bfseries\lr{Infants of DM mothers} ـ ۱۴۷}\par
\lr{Subject: Nelson 22}
\end{center}

\IDMpage{۱}
\begin{itemize}
\item عوارض جنینی دیابت بارداری (\lr{GDM}) کمتر از دیابت قبل از بارداری است.
\item اگر مادر دیابت پیش از بارداری داشته باشد، باید بارداری خواسته باشد.
  \begin{itemize}
  \item \lr{Tight Control}: \IDMmath{\mathrm{HbA}_{1c}<6.5\%}.
  \item اگر \IDMmath{\mathrm{HbA}_{1c}>7\%}: افزایش ریسک امبریوپاتی دیابتی (\lr{DEP}).
  \item نقص لولهٔ عصبی (\lr{NTD})، نقص قلبی، \lr{Caudal regression syndrome}.
  \item \IDMmath{\uparrow} خطر سقط خودبه‌خودی.
  \end{itemize}
\item در \lr{GDM}، بروز هایپرگلاسمی در بعد از هفتهٔ ۲۵ بارداری رخ می‌دهد؛ خطر مالفورمیشن جنینی کمتر است.
\item ۲۵ تا ۴۵\% شیرخواران (\lr{IDM}) مبتلا به \lr{LGA} هستند.
\item محدودیت رشد جنین در مادر با دیابت قبلی که عارضهٔ عروقی دارد دیده می‌شود ولی شیوع آن خیلی خیلی کمتر است.
\item یکی از علل \lr{Fetal loss} طی بارداری در مادر دیابتی، \lr{DKA} است.
\item خطر مرگ نوزاد از مادر دیابتی \IDMunclear{۵} برابر است.
\end{itemize}
\subsection*{پاتوفیزیولوژی}
یک تئوری این است که هایپرگلاسمی مادر منجر به هایپرگلاسمی جنین می‌شود؛ پاسخ جنین: هایپرانسولینیسم جنین.

\clearpage
\IDMpage{۲}
\begin{itemize}
\item با اینکه گلوکز مادر از جفت عبور می‌کند، انسولین مادر و \IDMunclear{گلوکاگون} از جفت عبور نمی‌کنند.
\item هایپرانسولینیسم و هایپرگلاسمی جنین: هایپرانسولینیسم جنین.
  \begin{itemize}
  \item افزایش \lr{uptake} کبدی گلوکز و سنتز گلیکوژن، افزایش لیپوژنز.
  \item تقویت سنتز پروتئین.
  \item هایپرتروفی و هایپرپلازی سلول‌های جزایر پانکراس.
  \item افزایش وزن جفت و ارگان‌هأ جنین؛ بجز مغز.
  \item هایپرتروفی میوکارد.
  \item \IDMmath{\uparrow} اکسترامدولاری هماتوپوئزیس.
  \end{itemize}
\item هایپرانسولینیسم و هایپرگلاسمی: اسیدمی در جنین؛ افزایش خطر مرده‌زایی.
\item زمانی که ارتباط جفتی مادر با جنین (نوزاد) قطع می‌شود: خطر هایپوگلاسمی نوزاد در ساعت‌های اولیهٔ تولد.
\item باید گلوکز مادر طی \lr{labor} و \lr{delivery} به صورت \lr{Tight Control} باشد؛ خطر هایپوگلاسمی نوزاد کم می‌شود.
\item سطح انسولین در شیرخوار با مادر \lr{GDM} کمتر از مادر با دیابت قبل از بارداری است.
\item سطح انسولین \IDMunclear{جفتی} شیرخوار با مادر \lr{PGDM} بسیار بالاتر از شیرخوار عادی است، علی‌رغم اینکه سطح گلوکز هر دو یکسان باشد.
  \begin{itemize}
  \item به دنبال دریافت گلوکز، به طور غیرمتناسب سطح انسولین بالا ترشح می‌شود.
  \item دادن \lr{Arginine} پاسخ به انسولین را تسهیل می‌کند.
  \end{itemize}
\end{itemize}

\clearpage
\IDMpage{۳}
{\footnotesize یادداشت‌های بالای صفحه: \lr{NPO}؛ \IDMunclear{۲۰۷}؛ \IDMunclear{«هایپرتروفیک» و واژهٔ کادربندی‌شدهٔ روی‌هم‌نویسی‌شده در گوشهٔ چپ}.\par}
\begin{itemize}
\item علت اصلی هایپوگلاسمی نوزاد، هایپرانسولینیسم است ولی کاهش پاسخ به اپی‌نفرین و گلوکاگون نیز روند را تشدید می‌کند.
\item غربالگری و درمان تهاجمی هایپوگلاسمی نوزادی باید انجام شود.
\item نوزاد ماکروزوم در خطر افزایش‌یافتهٔ ترومای حین زایمان است (آسیب \lr{BNP}).
  \begin{itemize}
  \item خطر آسفیکسی: هم به علت سایز بالا، هم تحمل استرس پایین؛ خصوصاً اگر کاردیومیوپاتیک شوند.
  \end{itemize}
\item هایپوگلاسمی:
  \begin{itemize}
  \item در ۲۵ تا ۵۰\% نوزاد از مادر \lr{pGDM}.
  \item در ۱۵ تا ۲۵\% نوزادان از مادر \lr{GDM}.
  \item تنها درصد کمی از نوزادان علامت‌دار می‌شوند.
  \end{itemize}
\item کمترین حد گلوکز نوزادی (\lr{nadir}) در ساعات ۱ تا ۳ تولد رخ می‌دهد.
  \begin{itemize}
  \item هایپوگلاسمی می‌تواند تا ۷۲ ساعت و در موارد نادر تا ۷ روز باقی بماند.
  \item درمان شیر دادن \lr{Frequent} است.
  \item برخی نوزادان به دکستروز \lr{IV} نیاز پیدا می‌کنند؛ علائم نوزاد.
  \end{itemize}
\end{itemize}
{\footnotesize یادداشت جدا در حاشیهٔ چپ: \IDMunclear{\lr{N.T b-ll}}.\par}
\subsection*{تظاهرات بالینی}
\begin{itemize}
\item ماکروزومی (به دلیل افزایش \lr{body fat}).
\item صورت پف‌آلود و پلوتوریک مانند کوشینگ.
\end{itemize}

\clearpage
\IDMpage{۴}
\subsection*{نوزاد دچار هایپوگلاسمی}
\begin{itemize}
\item \lr{Jittery}، ترموروس، هایپراکسیتابل طی ۳ روز اول.
\item هایپوتونی، لتارژی، \lr{poor feeding} و شاید تشنج.
\item اگر این علائم زودتر ایجاد شوند، به نفع هایپوگلاسمی هستند ولی می‌توانند به علت \IDMmath{\mathrm{Ca}\downarrow} و \IDMmath{\mathrm{Mg}\downarrow} نیز باشند.
  \begin{itemize}
  \item طی ۲۴ تا ۷۲ ساعت اول رخ می‌دهند؛ به علت تأخیر در پاسخ پاراتیروئید هستند.
  \end{itemize}
\item آسفیکسی و تحریک‌پذیری سیستم عصبی مرکزی، ریسک هایپوگلاسمی، \IDMmath{\mathrm{Mg}\downarrow} و \IDMmath{\mathrm{Ca}\downarrow} را نیز افزایش می‌دهد.
\end{itemize}
\subsection*{تاکی‌پنه و عوارض تنفسی}
\begin{itemize}
\item تاکی‌پنه در بسیاری از \lr{IDMs} طی ۲ روز نخست رخ می‌دهد.
\item می‌تواند ثانویه به هایپوگلاسمی، هایپوترمی، پلی‌سایتمی، نارسایی قلبی، \lr{TTN} یا ادم مغزی باشد.
\item شیوع \lr{RDS} در \lr{IDMs} بالاتر است.
  \begin{itemize}
  \item اثر مهاری انسولین بر بیان پروتئین سورفاکتانت.
  \end{itemize}
\item هیپوپلازی پارانشیم به همراه \lr{RDS} رخ می‌دهد که هر دو ثانویه به هایپرانسولینیسم هستند.
\end{itemize}
\subsection*{کاردیومگالی}
\begin{itemize}
\item شایع است. نارسایی قلبی در ۵ تا ۱۰\% موارد \lr{IDMs} رخ می‌دهد.
\item به دنبال هایپرتروفی سپتوم بین‌بطنی.
\end{itemize}
\begin{latin}
Transient Idiopathic Hypertrophic Subaortic Stenosis
\end{latin}

\clearpage
\IDMpage{۵}
\subsection*{ادامهٔ کاردیومگالی}
\begin{itemize}
\item داروهای اینوتروپ، انسداد را تشدید می‌کنند؛ کنترااندیکاسیون.
\item بتابلاکر انسداد را رفع می‌کند.
\item این وضعیت به طور خودبه‌خودی و تدریجی مرتفع می‌شود.
\end{itemize}
\subsection*{اختلالات نورولوژیک حاد}
\begin{itemize}
\item لتارژی، تحریک‌پذیری، \lr{poor feeding}.
\item اگر به علت هایپوگلاسمی باشد، با رفع آن علائم برطرف می‌شوند.
\item اگر به علت آسفیکسی باشد: تا هفته‌ها باقی می‌مانند.
\end{itemize}
\begin{itemize}
\item ریسک هایپربیلی‌روبینمی، پلی‌سایتمی، فقر آهن، ترومبوز ورید کلیوی در \lr{IDMs} بالاتر است.
\item هم‌زمان شک به ترومبوز ورید کلیوی: تودهٔ فلانک، هماچوری، ترومبوسایتوپنی.
\item ریسک آنومالی‌های مادرزادی در \lr{IDMs} چهار برابر است؛ با حد \lr{HbA1c} طی ۳ ماههٔ اول که ارگانوژنز رخ می‌دهد بالا می‌رود.
\end{itemize}
\subsection*{آنومالی‌های شایع}
\begin{itemize}
\item عدم بسته شدن نورال تیوب (آنانسفالی، مننگومیلوسل، انسفالوسل).
\item \lr{TGV} (جابجایی عروق بزرگ)، \lr{ASD}، \lr{VSD}.
\item استنوز آئورت (\lr{AS})، کوآرکتاسیون آئورت.
\item سایر: آترزی روده، سندرم کودال آژنزیس، آژنزی کلیه، هیدرویورتر، کلیهٔ نعل \IDMunclear{اسب؛ واژهٔ پایانی کوتاه شده است}.
\end{itemize}

\clearpage
\IDMpage{۶}
\begin{itemize}
\item سندرم کولون چپ کوچک: یک آنومالی نادر؛ در سه‌ماههٔ دوم و سوم رخ می‌دهد.
\end{itemize}
\subsection*{درمان}
\begin{itemize}
\item کنترل \lr{DM} در مادر.
\item \lr{NST} جنین (\lr{nonstress test}) + پروفایل بیوفیزیکال.
\item \lr{Doppler Velocimetry}.
\item زایمان حتماً در بیمارستان باشد.
\item زنان \lr{DM-1}: لزوم \lr{tight Cont} طی بارداری:
  \begin{itemize}
  \item \IDMmath{\mathrm{HbA}_{1c}<6.5\%}.
  \item \lr{Average Daily} \IDMmath{<95\,\mathrm{mg/ml}} (واحد مطابق دست‌نویس).
  \item بدون هایپوگلاسمی.
  \end{itemize}
\item درمان \lr{GDM}: رژیم غذایی، مانیتورینگ قند خون، انسولین، داروهای آنتی‌هایپرگلاسمیک خوراکی: متفورمین، \lr{Glyburide}.
\item فارغ از اندازه، بدو تولد نوزاد \lr{IDM}، همه باید \lr{Close Observation}.
\item تست غربالگری قند خون نوزاد باید بین \lr{30 min} و \lr{1 h} انجام شود.
\item \lr{AAP} گایدلاین: سطح گلوکز حداقل ۴۰ میلی‌گرم برای ۴ ساعت اول حفظ شود؛ حداقل ۴۵ برای ساعت ۴ تا ۲۴.
\item \lr{PES} گایدلاین: طی ۴۸ ساعت اول بالای ۵۰ میلی‌گرم باشد. (عدد قبلی خط‌خورده و ۵۰ جایگزین شده است.)
\item خط اول، فرم بدون علامت، درمان: هایپوگلاسمی بدون علامت، \lr{Feeding}؛ ژل دکستروز، «کاملاً ایمن».
\end{itemize}
{\footnotesize یادداشت‌های پایین صفحه: \IDMunclear{\lr{D. 298}}، \IDMunclear{۱۰۰٫۱۰}.\par}

\clearpage
\IDMpage{۷}
\begin{itemize}
\item تغذیهٔ نوزاد ظرف ۱ ساعت اولیه آغاز گردد.
\end{itemize}
\subsection*{اگر گلوکز را چک کردیم}
% Inequalities reproduce the handwritten branches, including the strict <49.
\begin{itemize}
\item \IDMmath{\geq 50}: هشت تا دوازده مرتبه \lr{feeding} در ۲۴ ساعت.
\item \IDMmath{30\leq\alpha<49}: ژل دکستروز ۴۰\%؛ چک گلوکز ۳۰ دقیقه بعد.
  \begin{itemize}
  \item چک مجدد، \IDMmath{\geq 50}: مشکلی ندارد.
  \item \IDMmath{30\leq\alpha<49}: گلوکز ۴۰\% ژل یا وریدی.
  \item \IDMmath{<30}: در \lr{NICU} دکستروز بگیرد.
  \end{itemize}
\item \IDMmath{<30}: دکستروز وریدی + ژل گلوکز بوکوس \IDMunclear{بوکوس/بوکال}.
\end{itemize}
\begin{itemize}
\item دوز وریدی مداوم دکستروز: \IDMmath{4\text{--}8\,\mathrm{mg/kg/min}}.
\item دوز بولوس \lr{D/W 25\%}: هایپرتونیک است؛ کنترااندیکاسیون مطلق.
\item شواهد نورولوژیک هایپوگلاسمی، عدم تحمل \lr{PO}: \lr{D/W} وریدی.
\end{itemize}

\clearpage
\section*{تصاویر اصلی}
{\footnotesize در هر هفت تصویر، سربرگ چاپی \lr{Subject:} و \lr{Date: / /} و نشان چاپی \lr{YASHA} وجود دارد. شمارهٔ مبحث ۱۴۷ است. شماره‌های دست‌نویس پایین صفحات و خط‌خوردگی‌ها در تصاویر زیر حفظ شده‌اند.\par}
\IfFileExists{images/pic01.jpg}{\def\IDMimagebase{images/}}{%
  \def\IDMimagebase{pediatrics/ped-neonatology/infants of DM mothers/images/}}
\begin{center}
\includegraphics[width=\linewidth,height=0.84\textheight,keepaspectratio]{\IDMimagebase pic01.jpg}
\end{center}
\clearpage
\begin{center}
\includegraphics[width=\linewidth,height=0.95\textheight,keepaspectratio]{\IDMimagebase pic02.jpg}
\end{center}
\clearpage
\begin{center}
\includegraphics[width=\linewidth,height=0.95\textheight,keepaspectratio]{\IDMimagebase pic03.jpg}
\end{center}
\clearpage
\begin{center}
\includegraphics[width=\linewidth,height=0.95\textheight,keepaspectratio]{\IDMimagebase pic04.jpg}
\end{center}
\clearpage
\begin{center}
\includegraphics[width=\linewidth,height=0.95\textheight,keepaspectratio]{\IDMimagebase pic05.jpg}
\end{center}
\clearpage
\begin{center}
\includegraphics[width=\linewidth,height=0.95\textheight,keepaspectratio]{\IDMimagebase pic06.jpg}
\end{center}
\clearpage
\begin{center}
\includegraphics[width=\linewidth,height=0.95\textheight,keepaspectratio]{\IDMimagebase pic07.jpg}
\end{center}
\IDMfinish
